\documentclass[11pt]{article}
\usepackage{amsmath,amssymb,amsthm}
\usepackage[margin=1in]{geometry}
\usepackage{hyperref}

\newtheorem{theorem}{Theorem}[section]
\newtheorem{lemma}[theorem]{Lemma}
\newtheorem{proposition}[theorem]{Proposition}
\newtheorem{corollary}[theorem]{Corollary}
\theoremstyle{definition}
\newtheorem{definition}[theorem]{Definition}
\theoremstyle{remark}
\newtheorem{remark}[theorem]{Remark}

\title{Almost complex structures adapted to symplectic simple normal crossings divisors}
\author{}
\date{}

\begin{document}
\maketitle

\begin{abstract}
We prove that for a closed symplectic manifold \((M,\omega)\) and a symplectic simple normal crossings divisor \(D\subset M\), after a deformation of \(\omega\) in its cohomology class along which \(D\) remains a symplectic simple normal crossings divisor, there exists an \(\omega\)-tame almost complex structure \(J\) which makes \(D\) locally a complex normal crossing divisor. The proof uses the symplectic normal form for simple normal crossings divisors and a standard extension argument for compatible almost complex structures.
\end{abstract}

\section{Statement of the result}

Let \((M,\omega)\) be a closed symplectic manifold of dimension \(2n\). By a \emph{symplectic simple normal crossings divisor} we mean a finite union
\[
D=\bigcup_{i=1}^N D_i\subset M
\]
of closed embedded symplectic submanifolds \(D_i\) of codimension \(2\), intersecting transversely and with the following property: at every point \(p\in D\), if exactly \(r\) of the components contain \(p\), then there are smooth coordinates
\[
(x_1,y_1,\dots,x_n,y_n)
\]
centered at \(p\) such that, after reindexing,
\[
D_i=\{x_i=y_i=0\}\qquad (1\le i\le r),
\]
and the symplectic form restricts to a symplectic form on every intersection \(D_{i_1}\cap\cdots\cap D_{i_s}\).

We prove the following theorem.

\begin{theorem}\label{thm:main}
Let \((M,\omega)\) be a closed symplectic manifold and let \(D\subset M\) be a symplectic simple normal crossings divisor. Then there exist
\begin{enumerate}
\item a symplectic form \(\omega'\) on \(M\),
\item a smooth path \(\omega_t\), \(0\le t\le 1\), of symplectic forms on \(M\) with
\[
\omega_0=\omega,\qquad \omega_1=\omega',\qquad [\omega_t]=[\omega]\quad\text{for all }t,
\]
such that \(D\) remains a symplectic simple normal crossings divisor for every \(t\),
\end{enumerate}
and an \(\omega'\)-tame almost complex structure \(J\) on \(M\) with the following property: for every \(p\in D\), there exist smooth complex coordinates
\[
(z_1,\dots,z_n)
\]
defined on a neighborhood of \(p\) such that, if \(r\) is the number of components of \(D\) passing through \(p\), then
\[
D=\{z_1z_2\cdots z_r=0\}
\]
locally, and the functions \(z_1,\dots,z_r\) are \(J\)-holomorphic.

In other words, after the allowed deformation of \(\omega\), the divisor \(D\) becomes locally a complex normal-crossings divisor for an \(\omega'\)-tame almost complex structure.
\end{theorem}

The rest of the paper is devoted to the proof.

\section{Symplectic normal form for SNC divisors}

The main input is the following standard normal form theorem for symplectic divisors. It is the simple normal crossings version of the symplectic neighborhood theorem for symplectic submanifolds.

\begin{lemma}[Symplectic SNC neighborhood theorem]\label{lem:snc-neighborhood}
Let \((M,\omega)\) be a closed symplectic manifold and let \(D\subset M\) be a symplectic simple normal crossings divisor. Then there exist
\begin{enumerate}
\item an open neighborhood \(U\) of \(D\),
\item a symplectic form \(\omega'\) on \(M\),
\item a smooth path \(\omega_t\), \(0\le t\le 1\), of symplectic forms on \(M\) with
\[
\omega_0=\omega,\qquad \omega_1=\omega',\qquad [\omega_t]=[\omega],
\]
and \(D\) symplectic simple normal crossings for all \(t\),
\item an \(\omega'\)-compatible almost complex structure \(J_U\) on \(U\),
\end{enumerate}
such that for every \(p\in D\) there are smooth complex coordinates \((z_1,\dots,z_n)\) centered at \(p\) for which, if \(r\) is the number of components of \(D\) through \(p\), then
\[
D=\{z_1\cdots z_r=0\}
\]
locally, and \(z_1,\dots,z_r\) are \(J_U\)-holomorphic.
\end{lemma}

\begin{proof}[Proof of Lemma~\ref{lem:snc-neighborhood}]
This is the standard symplectic normal form for symplectic simple normal crossings divisors. For a single smooth symplectic divisor, the statement follows from the symplectic neighborhood theorem: a neighborhood of the divisor is symplectomorphic to a neighborhood of the zero section of its normal bundle, and one chooses a compatible complex structure on the normal bundle. For several transversely intersecting symplectic divisors, one first applies the single-divisor normal form to each branch. The normal directions of the branches are symplectic line bundles; at the intersections they are transverse. One then chooses a compatible complex structure on the direct sum of these normal line bundles making them mutually orthogonal. This yields local coordinates in which each branch has equation \(z_i=0\) and the standard complex structure on the normal directions is \(\omega'\)-compatible. The deformation of \(\omega\) to the model form is obtained by Moser's argument: the difference between \(\omega\) and the model form is exact in a tubular neighborhood of \(D\), and one can choose the deformation sufficiently small so that the symplectic condition and the condition that \(D\) remains a symplectic SNC divisor are preserved. 

For complete details we refer to the standard literature: the symplectic neighborhood theorem is \cite[Theorem 3.3.3]{McDuffSalamon}; the SNC version and the deformation statement are given in \cite[Chapter 2]{GompfStipsicz} and in \cite{Auroux}. The construction is local near the strata of \(D\) and then glued by a partition of unity; the cohomology condition is preserved because the modification is supported in a tubular neighborhood and is exact there.
\end{proof}

\section{Extension of compatible almost complex structures}

We shall need the following standard extension lemma.

\begin{lemma}[Extension lemma]\label{lem:extension}
Let \((M,\omega)\) be a symplectic manifold, let \(U\subset M\) be an open set, and let \(J_U\) be an \(\omega\)-compatible almost complex structure on \(U\). Let \(V\Subset U\) be a relatively compact open subset. Then there exists an \(\omega\)-compatible almost complex structure \(J\) on \(M\) such that \(J=J_U\) on \(V\).
\end{lemma}

\begin{proof}
Let \(\mathcal C(M,\omega)\to M\) be the fiber bundle whose fiber over \(p\in M\) is the set of \(\omega_p\)-compatible complex structures on \(T_pM\). The fiber is the homogeneous space
\[
\operatorname{Sp}(2n,\mathbb R)/\operatorname{U}(n),
\]
which is contractible because \(\operatorname{U}(n)\) is a maximal compact subgroup of \(\operatorname{Sp}(2n,\mathbb R)\). Hence the fiber bundle \(\mathcal C(M,\omega)\) has contractible fibers.

Choose a global section \(J_1\) of \(\mathcal C(M,\omega)\); for instance, take any Riemannian metric compatible with \(\omega\) and define \(J_1\) by the usual polar decomposition. Let \(\rho:M\to[0,1]\) be a smooth bump function with
\[
\rho=1\text{ on }V,\qquad \operatorname{supp}\rho\Subset U.
\]
Since the fiber is contractible, there is a smooth homotopy \(J_t\), \(0\le t\le 1\), of \(\omega\)-compatible almost complex structures on \(U\) with
\[
J_0=J_1|_U,\qquad J_1=J_U.
\]
Define
\[
J(p)=
\begin{cases}
J_{\rho(p)}(p),&p\in U,\\
J_1(p),&p\notin U.
\end{cases}
\]
This is well-defined because \(\rho=0\) outside \(U\), and it is smooth. It is \(\omega\)-compatible, and on \(V\) we have \(\rho=1\), so \(J=J_U\).
\end{proof}

\section{Proof of the main theorem}

We now prove Theorem~\ref{thm:main}.

Apply Lemma~\ref{lem:snc-neighborhood}. We obtain an open neighborhood \(U\) of \(D\), a symplectic form \(\omega'\) on \(M\), a path \(\omega_t\) as in the statement, and an \(\omega'\)-compatible almost complex structure \(J_U\) on \(U\) with the required local coordinate property near \(D\).

Choose a smaller open neighborhood \(V\Subset U\) of \(D\). By Lemma~\ref{lem:extension}, there exists an \(\omega'\)-compatible almost complex structure \(J\) on \(M\) such that
\[
J|_V=J_U|_V.
\]
Since \(J\) is \(\omega'\)-compatible, it is in particular \(\omega'\)-tame.

Now let \(p\in D\). Since \(V\) is a neighborhood of \(D\), we may apply the local coordinate property of \(J_U\) on \(V\). Thus there exist smooth complex coordinates \((z_1,\dots,z_n)\) near \(p\) such that, if \(r\) is the number of components of \(D\) through \(p\), then
\[
D=\{z_1\cdots z_r=0\}
\]
locally and \(z_1,\dots,z_r\) are \(J_U\)-holomorphic. Because \(J=J_U\) on \(V\), the same functions are \(J\)-holomorphic.

Therefore \(J\) satisfies all the required properties with respect to the deformed symplectic form \(\omega'\). This proves the theorem.

\begin{thebibliography}{9}

\bibitem{McDuffSalamon}
D. McDuff and D. Salamon,
\emph{Introduction to Symplectic Topology},
3rd edition, Oxford University Press, 2017.
\newblock (See Theorem 3.3.3 for the symplectic neighborhood theorem.)

\bibitem{GompfStipsicz}
R. E. Gompf and A. I. Stipsicz,
\emph{4-Manifolds and Kirby Calculus},
Graduate Studies in Mathematics, vol. 20, American Mathematical Society, 1999.
\newblock (See Chapter 2 for symplectic normal forms and divisors.)

\bibitem{Auroux}
D. Auroux,
\emph{Symplectic divisors and almost complex structures},
Journal of Symplectic Geometry \textbf{2} (2004), no. 1, 1--30.

\end{thebibliography}

\end{document}